\documentclass[10pt,a4paper]{amsart}
\usepackage[margin=19mm]{geometry}
\usepackage{amsmath,amssymb,mathtools}
\usepackage[scr=rsfs,cal=euler]{mathalfa}
\usepackage{xfrac}
\usepackage[unicode,hidelinks,bookmarks=false]{hyperref}
\newcommand{\ie}{i.e.\ }
\newcommand{\cf}{cf.\ }
\newcommand{\resp}{respectively\ }
\newcommand{\Subsection}{\subsection}
\providecommand{\nobreakdash}{\nobreak\hbox{-}}
\setlength{\emergencystretch}{2em}
\multlinegap0pt
\usepackage{amssymb}
\usepackage[normalem]{ulem}
\newtheorem{theorem}{Theorem}
\newcommand{\Si}{\Sigma}
\renewcommand{\a}{\alpha}
\renewcommand{\b}{\beta}
\newcommand{\be}{\begin{equation}}
\newcommand{\cC}{{\mathcal C}}
\newcommand{\cD}{{\mathcal D}}
\newcommand{\cT}{{\mathcal T}}
\newcommand{\e}{\varepsilon}
\newcommand{\ee}{\end{equation}}
\newcommand{\g}{\gamma}
\renewcommand{\k}{\kappa}
\renewcommand{\l}{\lambda}
\newcommand{\p}{\partial}
\newcommand{\s}{\sigma}
\newcommand{\w}{\widetilde}
\title{Well-Posed Geometric Boundary data in General Relativity, III: \\ Conformal-mean curvature boundary data}
\author{Zhongshan An, Michael T. Anderson}
\date{Source: arXiv:2503.12599v3 (CC BY 4.0)}
\begin{document}
\maketitle

\noindent\emph{Source and licence.} Well-Posed Geometric Boundary data in General Relativity, III: \\ Conformal-mean curvature boundary data. Version arXiv:2503.12599v3 (CC BY 4.0), distributed by arXiv under the Creative Commons Attribution 4.0 International licence, http://creativecommons.org/licenses/by/4.0/, as recorded in the arXiv licence metadata for this exact version and shown on the abstract page https://arxiv.org/abs/2503.12599v3. Full licence notice: \texttt{licenses/arxiv-2503.12599-CC-BY-4.0.txt}.\par\smallskip

\noindent\textit{Editorial scope.} Counted unit: the article's theorem globexist (source0.tex lines 1531--1538). The article's complete original proof of it is delivered with it (source0.tex lines 1540--1580, one \texttt{proof} environment of 41 lines). No mathematics is written, repaired or re-derived: the statement and the proof are the article's own lines, in the article's own notation, and no retained line is substituted or re-flowed.  Retained uncounted as the source-local context the proof needs: lines 590--605; lines 1336--1356.  Nothing was omitted from any retained span, so the package carries no omission record.  No retained line contains a citation, so the wrapper carries no bibliography and no bibliography entry.  No retained line contains a figure environment, an image-inclusion command or drawing code, so no image is delivered and the package declares no asset dependency.  Editorial formatting only, in addition: an AMS \texttt{amsart} preamble that restates the article's own declarations and macros used by the retained lines, namely the environments \texttt{theorem} on separate counters, together with the standard packages those lines need.  The retained environments print in this document as Theorem 1 in order of appearance, whereas the article prints the same items with the numbering of its own full document.  The complete native source of the article as distributed in the arXiv e-print for this exact version is bundled unchanged as \texttt{sources/174865/source0.tex}.
  When $(M, g)$ is of size bounded away from $0$ and $\infty$, $U$ will be a small domain metrically, and correspondingly, the coordinates 
$x^{\a}$ will vary only over a small range. To renormalize this situation, as usual the metric and coordinates are rescaled simultaneously; thus 
for $\l$ small, set 
$$\w g = \l^{-2}g, \ \ \w x^\a=\l^{-1}x^\a$$
so that 
\be\label{rescale}
\w g(\p_{\w x^\a},\p_{\w x^\b})|_{\w x}=g(\p_{ x^\a},\p_{ x^\b})|_{\l x}. 
\ee
The equation \eqref{rescale} holds in the same way for any variation $h = \frac{d}{ds}(g + sh)|_{s=0}$ of $g$. Note that while the components 
$g_{\a\b}$ of $g$ are invariant under such a rescaling, all higher derivatives become small: 
$$\p_{\w x^{\mu}}^k \w g_{\a\b} = \l^k \p_{x^{\mu}}^k g_{\a\b}.$$ 
Thus the coefficients are close to constant functions in the rescaled chart, 
\be \label{eps}
||\w g - g_{\a_0}||_{C^{\infty}(U)} \leq \e = \e(\l, g),
\ee
where $g_{\a_0}$ is a flat (constant coefficent) Minkowski metric. This is the frozen coefficient approximation. 

\begin{theorem}\label{locexist}
There exists $\e > 0$ such that, for any smooth metric $g \in \cD$ on $\w U$ which is $C^{\infty}$ $\e$-close to a standard Minkowski corner metric $g_{\a_0}$ 
as in \eqref{eps}, the linearization $D\Phi_g^H$ at $g$ satisfies: given any $C^{\infty}$ target data 
$\tau' \in \cD$ on $\w U$, there exists a $C^{\infty}$ variation $h \in T_g Met(\w U)$, such that 
\be \label{DwPhisol}
D\Phi_g^H(h) = \tau'.
\ee
Thus, the equation 
\be \label{LhF}
L(h) = F,
\ee
has a $C^{\infty}$ solution $h$ such that 
\be \label{IC1}
(g_S, K)'_h = (\g_S', \k'), \ \ (\nu_S)'_h = \nu_S', \ \ V'_h = V_S' \ \mbox{ on }S,
\ee
and
\be \label{BC1}
(([g_\cC]), H)'_h = (\s', \ell'), \ \ V'_h = V_\cC' \ \mbox{ on }\cC.
\ee
Further the constructed solution $h$ depends smoothly on the data $\tau'$. 
\end{theorem}

\begin{theorem}\label{globexist}
Let $(M, g)$, $g \in \cD$, be a smooth globally hyperbolic Lorentzian manifold, $M \simeq I \times S$, with compact Cauchy surface $S$ having compact 
boundary $\p S = \Si$ and with timelike boundary $(\cC, g_{\cC})$. Then for any smooth $\tau' \in \cD \subset T_{\tau}(\cT^H)$, the equation 
\be \label{wL3}
D\Phi_g^H(h) = \tau'
\ee
has a $C^{\infty}$ solution $h$. The constructed solution $h$ depends smoothly on the data $\tau'$. 
\end{theorem} 

 \begin{proof}

  Given the previous local results, the proof of Theorem \ref{globexist} is essentially standard. We provide the details below for completeness. 
    
  Given $(M, g)$ as in \S 1, first choose an open cover $\{U_i\}_{i=1}^N$ of the corner $\Si$, where each $U_i$ is small enough so that the local 
existence result Theorem \ref{locexist} holds for each $U_i$. Additionally, choose an open set $U_0$ in the interior $M\setminus \cC$ such that 
$(\cup_{i=1}^N U_i)\cup U_0$ also covers a tubular neighborhood of the initial surface, i.e. $S\times [0,t']$ for some time $t'>0$.

Let $\{\w U_i\}_{i=0}^N$ be a thickening of the cover $\{U_i\}_{i=0}^N$, so that $U_i \subset \w U_i$, as described in \S 2.1. Let 
$\rho_i$ be a partition of unity subordinate to the cover $\{\w U_i\}_{i=0}^N$, so that ${\rm supp}\rho_i\subset \w U_i$, $\rho_i = 1$ on 
$U_i$ and  $\sum_i \rho_i=1$. Let 
$$\tau' =\big(F, (\g', \k', \nu', V_S'), (\s', \ell',V_{\cC}'), a' \big),$$
be arbitrary $C^{\infty}$ data in $T(\cT^H)$ on $M$. Then the data $\rho_i \tau'$ has compact support in the sense of \S 2.1 in $\w U_i$ 
and by Theorem \ref{locexist} there exists a solution $h_i$ in $\w U_i$ satisfying 
$$D\Phi^H(h_i) = \rho_i \tau'.$$ 
Moreover, as noted at the end of \S 2.1, by the finite propagation speed property, the solution $h_i$ has compact support in the sense of \S 2.1 for 
$t \in [0,t_i]$ for some $t_i > 0$. Thus, $h_i$ extends smoothly as the zero solution on $M_{t_i} \setminus \w U_i$. 

   Similarly, in the interior region $U_0 \subset \w U_0$, let $h_0$ be the solution to the Cauchy problem 
$$L(h_0) = \rho_0 F \ \mbox{ in } \w U_0, \ \ (g_S, K_S, \nu_S, V_g)'_h = \rho_0(\g', \k', \nu', V_S') \ \mbox{ on }S_0.$$
Then again $h_0$ has compact support in $\w U_0$ for some time interval $[0, t_0]$ and as above extends to $M_{t_0}$. We relabel so that 
$t_0$ is a common time interval for all $h_i$, $i \geq 0$. 

  The sum 
$$h = \sum_{i=0}^N h_i,$$
is thus well-defined on $M_{t_0}$ and by linearity
$$D\Phi^H(h) = \tau'.$$ 
By the proof of Theorem \ref{locexist}, the constructed $h$ satisfies the estimate
\be \label{enest2}
||h||_{C^{m}(S_t)} \leq K,
\ee
for all $m$ and for $t \leq t_0$. Here $K$ is a constant $K = K(t_0, \tau')$, which depends only on the background metric $g$ and the 
target data $\tau' \in \cD$. 

  One may now continue the solution $h$ past the time $t_0$ to all $t < \infty$ in the usual way, assuming of course that the background metric 
$g$ satisfies this condition. Briefly, let $I_{t_0}$ denote the target initial data of $h$ at the Cauchy slice $\{t=t_0\}$ and form the target data 
$\tau'$ by replacing the initial data in $\tau'$ with $I_{t_0}$. Then applying Theorem \ref{locexist} with initial slice $S_{t_0}$ and using \eqref{enest2} 
gives the continuation of $h$ to $[0, t_1]$, with $t_1 > t_0 + \mu$, where $\mu$ depends only on the boundary data $b'_h$ on $\cC$. 
Assuming then $g$ is defined for all time $t$, since then $b'_h$ is also globally defined on $\cC$ for all $t$, the same holds for $h$. 
  
\end{proof}

\end{document}
